\chapter{Conclusions}
\label{chap:conclusions}

\begin{introduction}
This chapter revisits the research questions in light of the work presented, summarises the contributions, honestly sets out the limitations of the current system, and describes the future work that follows most directly from it, including the viewer's intended release as open-source tooling for other researchers.
\end{introduction}

\section{Revisiting the Research Questions}
\label{sec:conc-rqs}

\paragraph{RQ1: a single pattern for arbitrary entity types.} This is answered affirmatively. Chapter~\ref{chap:architecture} presented an entity model built on one registration table and two escalating extension points, and Chapter~\ref{chap:implementation} showed that roughly thirty entity types, including static infrastructure, telematics equipment, simulated vehicles, and fire simulations, are supported through it, with all but one requiring no bespoke code in the generic spawn, update, or rendering paths. The worked example in Section~\ref{sec:impl-worked-example} showed that adding a further type is a two- or three-step change confined to a schema declaration, one table line, and an optional presentation class. The most demanding case, fire, fits the deepest extension point without forcing an exception.

\paragraph{RQ2: backend independence.} This is answered both by construction and by an actual interface. Backend knowledge is confined to the service layer described in Chapter~\ref{chap:architecture}, verifiable by inspection since no code above it names an HTTP verb, a URL, an authentication scheme, or the live-event envelope format; and, as Section~\ref{sec:impl-backend-interface} shows, that confined knowledge sits behind three small interfaces (for authentication and record access, for the live-connection protocol, and for decoding live-update messages) rather than only inside concrete classes. Pointing the viewer at a different backend is therefore writing one implementation of each interface and using it in place of the Ditto one, not rewriting the backend connection, the live-update stream, or the update router; describing that backend's entities is separate, ordinary work, the same work every new entity type already requires regardless of which backend supplies it (Section~\ref{sec:usecases}). No second implementation or offline mock has been written yet. Hence, the claim still rests on the existence and shape of the interfaces rather than a working substitution, but that remaining gap is a bounded implementation task, not an architectural one.

\paragraph{RQ3: scaling and rendering strategy.} This is answered. The ten-run sweep (Chapter~\ref{chap:evaluation}) rules out the naive \textsc{PerSign} strategy decisively: at 15,000 instances it issues 15--23\(\times\) more draw calls, spawns 37--38\(\times\) slower, and runs 50\% slower per frame than the two batched alternatives. Between those two, the result is narrower than the ``always atlas'' assumption this benchmark set out to test: \textsc{AtlasMaterial}'s much lower draw-call count makes it faster than \textsc{PerTexture} by less than 1\% of the frame time at the smaller counts, a statistically significant but practically negligible difference, and the two are not significantly different at 15,000 instances. At the same time, \textsc{PerTexture} wins decisively on memory (22\(\times\) less texture memory). The recommendation for this viewer's actual sign rendering is \textsc{PerTexture}, not the atlas approach the literature default would suggest.

\paragraph{RQ4: usability.} Answered positively, with qualifications. In a moderated study with ten participants from the intended user population (Section~\ref{sec:eval-user-testing}), the viewer scored a mean SUS of 81.1, between ``good'' and ``excellent'', and every participant completed the tasks covering camera control, entity selection, the detail panel, the geotechnical slope, the fire simulation, and navigation, with one performance-related failure on the ring-road task. The traffic task could not be run because the live traffic feed was unavailable, so that use case's usability is untested. The issues raised, chiefly keyboard-first navigation, poor entity visibility at low zoom, and missing feedback when placing an entity, belong to the camera and UI layers rather than to the decoupled design, and are addressable without architectural change.

\section{Summary of Contributions}
\label{sec:conc-contributions}

The four contributions stated in Chapter~\ref{chap:introduction} were delivered as follows. The first is an open-source digital twin viewer for mobility infrastructure, built on Unreal Engine and Cesium for Unreal, that renders roughly thirty entity types on a full-scale globe with live updates, with the entity model and backend-access layer isolated behind explicit seams. The second is the architectural pattern behind it: a service layer of small, single-purpose services that confines all backend contact, with the Ditto-specific pieces of that contact behind three small interfaces, and an entity model driven by a registration table with a stateless per-type extension point and a stateful per-instance component point, documented through the implementation decisions that realise it.

The third is a benchmark harness for the entity spawning and rendering path, comparing three mesh and material batching strategies for a high-cardinality type across a swept instance-count range, with each configuration averaged over ten runs. It rules out \textsc{PerSign} on draw calls, spawn time, and frame time, and recommends \textsc{PerTexture} over \textsc{AtlasMaterial} for this viewer's sign rendering on memory grounds, since the two are practically equivalent on frame time. The fourth is a user study with ten participants from the intended user population, evaluating the viewer through eight realistic tasks and the System Usability Scale.

\section{Limitations}
\label{sec:conc-limitations}

The main limitation is that backend decoupling has an interface but only one implementation. As set out under RQ2 and in Sections~\ref{sec:impl-backend-interface} and~\ref{sec:impl-replace-caveat}, the REST and authentication surface, the live-connection protocol, and the envelope format each sit behind their own interface. Still, exactly one implementation of each exists, no mock backend has been built, and the viewer has been validated against exactly one backend.

A second limitation is that the viewer inherits the backend's operational behaviour. Like any consumer application, it assumes the backend it points to behaves correctly; when the backend has connectivity or reliability problems of its own, the viewer is affected too, and no amount of code-level decoupling changes that.

\section{Future Work}
\label{sec:conc-future}

Future work falls into two groups: completing the architectural claims and making the viewer available beyond DT4MOB, and improving the application along the lines its users identified in the study in Section~\ref{sec:eval-user-testing}.

\subsection{Architecture and Release}

The most direct next step is to write a second implementation, or a mock, behind the existing interfaces. Section~\ref{sec:impl-backend-interface} describes the three backend interfaces, each with one Ditto implementation; writing a second implementation, or an offline mock that replays recorded traffic, would turn RQ2's architectural claim into a demonstrated one, and would let the viewer be developed and tested with no live backend at all. Beyond that, pointing the viewer at a different digital twin backend, or at a different Ditto deployment with a different entity schema, would be the strongest possible evidence for the decoupling.

A further step is to build an in-house, high-detail map of Portugal. Building the viewer's own map from high-detail geographic tiles, served as its own tile set, would give the project full control over the map's coverage, level of detail, and availability.

Finally, the viewer is intended for release as open-source software, so that other researchers and institutions working with digital twins of different infrastructure types can adopt it as a starting point rather than rebuilding a geospatial viewer from scratch. Such a release needs documentation of the extension points for external developers, a worked example repository adding a new entity type against a public dataset, a packaged build with the login flow pointed at a configurable backend, and a contribution guide. The architectural work in this thesis is what makes the release useful: a viewer that could only talk to one backend and show one project's entity types would not be worth adopting.

\subsection{Improvements from User Feedback}
\label{sec:conc-future-feedback}

The user study surfaced several concrete improvements. Almost all of them concern the camera and interface layers rather than the architecture, and most are inexpensive. They are discussed below by the part of the application they affect, in roughly the order of how often participants raised them.

The first concerns navigation. Keyboard-driven movement was the most frequent complaint, raised by roughly half the participants, who expected to pan and drag the globe with the mouse as in the map tools they use daily. The RTS camera should support click-and-drag panning (with the left or middle button) alongside the existing keys, so that the keyboard becomes optional rather than required. Participants also suggested rotation and vertical movement on familiar RTS-style keys, and one asked for camera movement to be clamped to the area of interest so the view cannot drift off the region with data. The FreeFly camera, in turn, was found too fast and visually heavy. It should offer an adjustable speed multiplier, reduce or make motion blur optional, and use Esc (rather than Tab) to release the cursor, which participants found more natural.

The second concerns finding entities at a distance. Participants found entities hard to see and to select when zoomed out, and hard to find without knowing in advance where they were. Three changes follow. Entity markers should grow with camera distance or have a minimum on-screen size, and use colour to distinguish types. Entities close together should collapse into a single cluster marker at far zoom levels, expanding as the camera approaches. A country-level overview should also mark where each digital twin's data is, so a user can jump to an area of interest instead of searching for it. To be consistent with the architecture, clustering and the overview should work on entity position alone, without per-type logic. A readout of the camera's current coordinates was also requested.

The third concerns feedback when placing entities and running simulations. Creating a fire simulation was the lowest-rated completed task, largely due to a lack of feedback. The placement flow should confirm that an entity was created, since one participant placed a second ignition point because they believed the first had failed. After \emph{Place Entity} is pressed, the cursor should change, or a prompt should appear, to show that a location is now expected; the entity type could also be chosen after pressing the button rather than before, which better matched participants' expectations. For the simulation itself, participants asked for a legend for the fire's colour scale and a progress indicator while it runs; one suggested a dedicated fire-simulation entry point rather than a generic placement tool.

The fourth concerns rendering quality. One participant's work machine could not run the viewer until a lighter build was made by changing engine rendering settings alone (disabling ray tracing and heavier anti-aliasing, among others), and some participants still saw slowdowns, notably after creating a fire simulation and when zooming in over the Porto region (Section~\ref{sec:eval-discussion}). Offering those settings as user-selectable quality presets, instead of fixing them at build time, would let one build serve both workstations and ordinary work machines.

Beyond these, one participant asked for dashboards that aggregate data across entities for analysis. Such dashboards already exist in Grafana, outside the viewer, and some entities link to them from their detail window; the request was to have them inside the viewer itself. Embedding them, or showing aggregated views alongside the map, would keep spatial inspection and analysis in one place.

Finally, the evaluation itself should be completed. The traffic task could not be run because the live feed was unavailable, so the traffic use case's usability is still untested; it should be re-run once the feed is back. A follow-up study after the changes above would also show whether they address the issues raised. It should record participants' hardware to separate interface problems from performance problems, and time each task, which this study did not.

\section{Closing Remarks}
\label{sec:conc-closing}

The central claim of this thesis is that a digital twin viewer for physical infrastructure need not be a closed, single-purpose application. By confining backend knowledge to a service layer and infrastructure-domain knowledge to a small set of opt-in extension points, the bulk of the viewer, meaning rendering, interaction, the camera, and coordinate placement, can be written once and left untouched as entity types are added and, in principle, as the backend changes. This is not free: the implementation in Chapter~\ref{chap:implementation} carries its own complexity, most visibly in the tile-streaming loop. But the result is a viewer whose cost of change is localised, and which is a plausible basis for reuse beyond the project it was built for.
